\documentclass[11pt,letterpaper]{article}
\usepackage[margin=0.9in]{geometry}
\usepackage[T1]{fontenc}
\usepackage{mathptmx,microtype,amsmath}
\usepackage[hidelinks]{hyperref}
\usepackage{titlesec}
\titleformat{\section}{\normalsize\bfseries}{\thesection}{0.6em}{}
\titlespacing*{\section}{0pt}{10pt}{5pt}
\setlength{\parindent}{0pt}
\setlength{\parskip}{6pt}
\pagestyle{plain}
\begin{document}
\begin{center}
{\Large\bfseries Scaling Laws for Biological Realism\\[3pt] in DNA Language Models}\par
\vspace{5pt}
{Kevin Song, Juan Ricardo}
\end{center}
\vspace{-5pt}

\section{Research Question and Motivation}
How does the length over which generated DNA preserves natural genomic patterns change with model size? We will compare different-sized DNA language models to test whether larger models maintain these patterns farther from an initial natural sequence and whether the improvement follows a predictable scaling law.

DNA language models can extend a short DNA sequence by repeatedly predicting and sampling the next base. Larger models generally predict natural DNA more accurately, but long generated sequences can still differ from real genomes. A sequence may look plausible near its starting point and gradually lose the patterns found in natural DNA. Understanding whether model size delays this deterioration would help determine what scaling achieves beyond better next-base prediction.

Here, realistic DNA means DNA that resembles natural sequences in measurable ways: the frequency of short base patterns, the lengths and spacing of repeats, and the arrangement of gene-coding regions. We will also test whether a classifier can tell generated DNA from natural DNA. These tests measure resemblance to a chosen organism's DNA; they do not establish that a sequence would function in a cell.

\section{Prior Work and What Remains Open}
Nguyen et al.\ introduced Evo and showed that DNA prediction improves with model scale across several architectures~\cite{evo}. Brixi et al.\ developed Evo 2, reporting improved prediction and generation across a broader range of organisms~\cite{evo2}. Both studies produced long DNA sequences with plausible genomic features, while identifying missing components in genome-scale generations.

Tzanakakis et al.\ found differences between generated and natural DNA in repeats, short sequence patterns, regulatory motifs, and organization over long distances~\cite{limitations}. Their accompanying software includes a convolutional neural network (CNN) that distinguishes natural from generated DNA and measures distinguishability at increasing distances from the initial prompt~\cite{benchmark}.

It remains unclear how the rate of deterioration changes with model size and whether different genomic patterns improve at different rates. We will compare models under matched generation settings to measure how long each preserves these patterns and test whether the relationship follows a consistent scaling law.

\newpage
\section{Plan}
\textbf{Train a controlled family of models.} We will use Evo 2's StripedHyena 2 architecture and its public training framework, Savanna~\cite{savanna}, to train models with approximately 5M, 10M, 25M, 50M, and 100M parameters. Each model will train from scratch on the same 100M DNA tokens, using the same data order, tokenizer, and context length. We will adjust model width and depth while keeping the mix of layer types as consistent as possible. Separate genomic regions will be reserved for validation and evaluation. This measures scaling at a fixed token budget; larger models will use more computation. Released 1B, 7B, and 40B Evo 2 base models will provide additional comparisons, with their different training histories reported separately.

\textbf{Measure how realism scales.} We will use identical natural prompts and sampling settings across models and examine equal-length windows at increasing distances from each prompt. The published CNN will measure how easily generated DNA can be distinguished from natural DNA. We will also develop a proposed \emph{structure retention score} based on repeat lengths and spacing. For each feature, we will measure the difference between generated and matched natural windows relative to the variation between natural windows. We will average these normalized differences to summarize structural change. We will test the score on natural DNA with controlled repeat changes and check whether it detects those changes when overall base composition is preserved. Separate genomic regions will be used to develop and test the score.

For each model, we will plot both measurements against distance and compare the rate of deterioration. We will test simple scaling relationships across the five small models and report variation across prompts and repeated generations. Prediction loss will be recorded alongside these results.

\textbf{Diagnose deterioration.} We will compare short continuations generated from natural prefixes with continuations from model-generated prefixes of the same length. Both groups will use the same amount of context and source organisms. If quality is worse after generated prefixes, this would support the explanation that the model's own history contributes to deterioration. We will also shorten the available context to test whether losing earlier information makes the problem worse. Feature-level results will show whether differences arise mainly in base composition, repeat lengths, or repeat spacing.

\textbf{Optional fine-tuning.} If time permits, we will explore reinforcement-learning fine-tuning of one small model. The model will sample short continuations, a frozen CNN will score them, and a policy-gradient update will increase the probability of higher-scoring outputs~\cite{seqgan}. We will retain next-base training on natural DNA and penalize large changes from the starting model to help keep training stable. We will compare this method with ordinary continued training under a similar computing budget. A separate classifier, the proposed structural score, and sequence-diversity measurements will test whether improvement extends beyond the reward used for training.

\section{First Step, Expected Results, and Alternatives}
The first step is to run a short training test of the 5M model in Savanna, verify next-base prediction and sampling, and estimate the cost of the full model family. We will then train the 5M and 10M models on 100M tokens each. An initial evaluation will use 100 held-out prompts, 256 bases per prompt, and 2,048-base continuations within a shared 4,096-base context limit. Successful runs will be followed by the remaining model sizes and longer continuations within that limit.

We expect larger models to preserve some genomic patterns for longer and hypothesize that generated prefixes will make deterioration worse. If the smallest models generate little useful structure, we will begin diagnosis with the 50M or 100M model. If the proposed score misses controlled repeat changes, we will revise its features before using it for model comparisons. If RL improves only its training reward, we will reduce the reward weight and examine which patterns the classifier favors. Scaling, metric development, and diagnosis will be the main focus; fine-tuning will be optional.

\clearpage
\begin{thebibliography}{9}
\small
\bibitem{evo} E. Nguyen et al. Sequence modeling and design from molecular to genome scale with Evo. \emph{Science}, 386:eado9336, 2024. \href{https://doi.org/10.1126/science.ado9336}{doi:10.1126/science.ado9336}.
\bibitem{evo2} G. Brixi et al. Genome modelling and design across all domains of life with Evo 2. \emph{Nature}, 652:1349--1361, 2026. \href{https://doi.org/10.1038/s41586-026-10176-5}{doi:10.1038/s41586-026-10176-5}.
\bibitem{limitations} A. Tzanakakis, I. Mouratidis, and I. Georgakopoulos-Soares. Fundamental limitations of genomic language models for realistic sequence generation. \emph{bioRxiv}, 2026. Preprint. \href{https://doi.org/10.64898/2026.01.17.700093}{doi:10.64898/2026.01.17.700093}.
\bibitem{benchmark} Georgakopoulos-Soares Lab. Synthetic Genomes Benchmark. Software accompanying Tzanakakis et al. Accessed October 17, 2026. \url{https://github.com/Georgakopoulos-Soares-lab/Synthetic_genomes_benchmark}.
\bibitem{savanna} Savanna contributors. Savanna: Pretraining infrastructure for multi-hybrid model architectures. Software. Accessed October 17, 2026. \url{https://github.com/Zymrael/savanna}.
\bibitem{seqgan} L. Yu, W. Zhang, J. Wang, and Y. Yu. SeqGAN: Sequence generative adversarial nets with policy gradient. \emph{AAAI}, 2017. \url{https://arxiv.org/abs/1609.05473}.
\end{thebibliography}
\end{document}
